\documentclass[11pt]{article}
\usepackage[T1]{fontenc}
\usepackage{lmodern}
\usepackage[margin=1in]{geometry}
\usepackage{amsmath,amssymb,amsthm,mathtools}
\usepackage{microtype}
\usepackage{enumitem}
\usepackage{xcolor}
\usepackage{tikz}
\usepackage[colorlinks=true,linkcolor=blue!45!black,citecolor=green!35!black,urlcolor=blue!55!black]{hyperref}
\usepackage[all]{hypcap}
\hypersetup{
  pdftitle={A product counterexample to the simplex maximum for projection-body volume},
  pdfsubject={An elementary counterexample in dimension twenty}
}
\newtheorem{theorem}{Theorem}
\newtheorem{lemma}[theorem]{Lemma}
\newtheorem{proposition}[theorem]{Proposition}
\newtheorem{corollary}[theorem]{Corollary}
\newcommand{\R}{\mathbb R}
\DeclareMathOperator{\conv}{conv}
\DeclareMathOperator{\vol}{vol}
\DeclareMathOperator{\proj}{proj}
\newcommand{\unverified}[1]{\par\smallskip\noindent\fbox{\parbox{0.95\linewidth}{\textbf{Unverified step.} #1}}\par\smallskip}
\newcommand{\step}[1]{\par\smallskip\noindent\textbf{#1}\ }
\title{A product counterexample to the simplex maximum\\
for projection-body volume}
\author{Rewritten for readability from the OpenAI manuscript of the same title}
\date{}
\raggedbottom

\begin{document}
\maketitle

\begin{abstract}
The product of two ten-dimensional simplices has larger normalized
projection-body volume than a twenty-dimensional simplex. This gives
a counterexample to Brannen's proposed simplex maximum. The proof is
elementary: a facet formula for projection bodies of polytopes, a
product identity, the exact simplex value, and one integer inequality.
\end{abstract}

\section{The result}\label{sec:result}

\subsection{Definitions}

Throughout, \(d\ge2\). A \emph{convex body} in \(\R^d\) is a compact
convex set with nonempty interior. We write \(|K|\) for the
\(d\)-dimensional volume of \(K\subset\R^d\), and \(\vol_j\) for
\(j\)-dimensional volume inside a \(j\)-dimensional affine subspace.
For a compact convex set \(C\), its \emph{support function} is
\[
 h_C(x)=\max_{y\in C}\langle x,y\rangle,\qquad x\in\R^d .
\]
\(S^{d-1}\) is the Euclidean unit sphere, \(\proj_{u^\perp}\) is the
orthogonal projection onto the hyperplane \(u^\perp\), \(B_2^j\) is the
Euclidean unit ball of \(\R^j\), and \(\kappa_j=|B_2^j|\).

The \emph{projection body} \(\Pi K\) (also \(\Pi_dK\)) of a convex body
\(K\subset\R^d\) is the compact convex set whose support function is
\[
 h_{\Pi K}(u)=\vol_{d-1}\bigl(\proj_{u^\perp}K\bigr)
 \quad(u\in S^{d-1}),
 \qquad
 h_{\Pi K}(tu)=t\,h_{\Pi K}(u)\quad(t\ge0).
\]
A compact convex set is determined by its support function (fact (S2)
in Section~\ref{sec:tools}). That a set with this support function
exists is proved in Lemma~\ref{lem:facets} for polytopes, which is all
the main theorem needs, and in Part~B of the proof of
Lemma~\ref{lem:affine} for general convex bodies.

The functional under study, and the constant it is compared with, are
\begin{equation}\label{eq:functional}
 R_d(K)=\frac{|\Pi K|}{|K|^{d-1}},
 \qquad c_d=\frac{(d+1)d^d}{d!}.
\end{equation}
Proposition~\ref{prop:simplex} shows that \(R_d(S)=c_d\) for every
full-dimensional simplex \(S\subset\R^d\). The statement ``simplices
maximize \(R_d\)'' is therefore the inequality
\begin{equation}\label{eq:proposed-bound}
 |\Pi K|\le c_d|K|^{d-1}\qquad\text{for all convex bodies }K\subset\R^d .
\end{equation}

\subsection{Main theorem}

\begin{theorem}\label{thm:counterexample}
Let \(T_{10}=\conv(0,e_1,\ldots,e_{10})\subset\R^{10}\), where the
\(e_i\) are the coordinate unit vectors, and let
\(K=T_{10}\times T_{10}\subset\R^{20}\). Then
\[
 \frac{R_{20}(K)}{c_{20}}
 =\frac{121\binom{20}{10}}{21\cdot2^{20}}
 =\frac{22\,355\,476}{22\,020\,096}>1.
\]
Thus \(K\) violates \eqref{eq:proposed-bound} in dimension twenty.
\end{theorem}

A quantitative extension, Corollary~\ref{cor:exponential-excess}, shows
that products of simplices beat the simplex value by a factor
\(\lambda^n\), with a fixed \(\lambda>1\), in every sufficiently large
dimension \(n\).

\subsection{Context}

The statements in this subsection are background quoted from the source
with its citations. None of them is used in a proof below.

Petty established the affine covariance of the projection-body
operator~\cite{Petty}; see also \cite[Equation~(1)]{Ludwig}. Brannen
conjectured in 1996 that simplices maximize \(R_d\) among all convex
bodies~\cite{Brannen}, which is \eqref{eq:proposed-bound}. It is the
volume of \(\Pi K\) itself that occurs here, not the volume of its
polar: the classical projection inequalities involving the polar concern
a different functional. For \(d\ge3\), Petty's separate minimum problem
asks whether ellipsoids are the only minimizers of the same functional
\(R_d\); see \cite[Section~1]{Saroglou} for these distinctions.

In dimension three, Chen, Feng, Li, Xi, and Xu settle the minimum
problem and also prove \(R_3(K)\le18=c_3\) for all three-dimensional
convex bodies, with tetrahedra attaining
equality~\cite[Theorems~1.1 and~1.2]{ChenEtAl}. Feng, Hu, Liu, and Xu
give counterexamples to \eqref{eq:proposed-bound} in every dimension
\(d\ge9\), using polytopes with at most \(d+2\)
facets~\cite[Theorem~1.3]{FengEtAl}; their construction uses perturbed
projections of a cube and Minkowski's existence
theorem~\cite[Section~4]{FengEtAl}. So the unrestricted-dimensional
simplex-maximum question already has a negative answer. The present
paper supplies a direct Cartesian-product construction with an exact
value. For centrally symmetric bodies, Henk gives the lower bound
\(2^d(9/8)^{\lfloor d/3\rfloor}\) for the optimal upper
constant~\cite{Henk}. The optimal upper constant in
\eqref{eq:proposed-bound} and its maximizing bodies remain separate
questions.

\subsection{Proof at a glance}

\textbf{Key idea.} The functional \(R_d\) is multiplicative under
Cartesian products of polytopes: \(R_{r+s}(A\times B)=R_r(A)R_s(B)\).
The simplex constants are not: \(c_{10}^2>c_{20}\). Hence
\(T_{10}\times T_{10}\) beats the \(20\)-simplex.

\textbf{Steps.}
\begin{enumerate}[label=\arabic*.,leftmargin=2em]
\item \emph{Facet formula} (Lemma~\ref{lem:facets}). For a polytope
 \(P\), \(\Pi P\) is a Minkowski sum of segments, one segment
 \([-\tfrac12s_F\nu_F,\tfrac12s_F\nu_F]\) per facet \(F\) (area
 \(s_F\), outward unit normal \(\nu_F\)).
\item \emph{Product identity} (Proposition~\ref{prop:product}). The
 facets of \(A\times B\) are \(F\times B\) and \(A\times G\). With
 step~1 this gives \(\Pi(A\times B)=(|B|\Pi A)\times(|A|\Pi B)\), and
 multiplicativity of \(R\) follows by counting exponents.
\item \emph{Affine invariance} (Lemma~\ref{lem:affine}).
 \(R_d(LP)=R_d(P)\) for invertible linear \(L\); translations also
 preserve \(R_d\). For polytopes this follows from step~1.
\item \emph{Simplex value} (Proposition~\ref{prop:simplex}). By step~1,
 \(\Pi T_d\) is a translate of \(([0,1]^d+[0,w])/(d-1)!\) with
 \(w=(1,\ldots,1)\). A Fubini argument gives
 \(|[0,1]^d+[0,w]|=d+1\), so \(R_d(T_d)=c_d\). Step~3 extends this to
 all simplices.
\item \emph{Arithmetic} (Section~\ref{sec:counterexample}). Steps~2
 and~4 give \(R_{20}(K)=c_{10}^2\), and \(c_{10}^2/c_{20}\) is the
 stated fraction.
\end{enumerate}

\textbf{Dependencies.} Lemma~\ref{lem:facets} feeds
Proposition~\ref{prop:product}, Part~A of Lemma~\ref{lem:affine}, and
Proposition~\ref{prop:simplex}. Lemma~\ref{lem:affine} is used only for
the last sentence of Proposition~\ref{prop:simplex} (all simplices have
the same value); its Part~B (non-polytopal bodies) is used nowhere
else. Theorem~\ref{thm:counterexample} and
Corollary~\ref{cor:exponential-excess} use only
Propositions~\ref{prop:product} and~\ref{prop:simplex} plus explicit
arithmetic.

\section{Conventions and tools}\label{sec:tools}

\textbf{Renamed symbols.} To keep one meaning per symbol, four letters of
the source are changed. In the proof of Lemma~\ref{lem:affine}, the
radius \(R\) becomes \(\rho\), the functions \(F,f_m\) become
\(g,g_m\), and the factors \(\lambda_m\) become \(\tau_m\). In
Corollary~\ref{cor:exponential-excess}, the constant \(A\) becomes
\(\Theta\), and the residue \(r\) with \(d_r=10+r\) becomes \(i\) with
\(q=10+i\).

\textbf{Minkowski sum.} \(C+D=\{c+d:c\in C,\ d\in D\}\). For
\(v\in\R^d\), \([a,b]\) denotes the segment with endpoints \(a,b\).

\textbf{Support-function facts.} Let \(C,D\) be compact convex sets.
\begin{enumerate}[label=(S\arabic*),leftmargin=3em]
\item \(h_{C+D}=h_C+h_D\); \(h_{tC}=t\,h_C\) for \(t\ge0\);
 \(h_{MC}(u)=h_C(M^{t}u)\) for a linear map \(M\); and for
 \(C\subset\R^r\), \(D\subset\R^s\),
 \(h_{C\times D}(u,v)=h_C(u)+h_D(v)\). Each is immediate from the
 definition: for instance
 \(\max_{c,d}\langle u,c+d\rangle=\max_c\langle u,c\rangle+\max_d\langle u,d\rangle\)
 and \(\langle u,My\rangle=\langle M^tu,y\rangle\).
\item If \(h_C\le h_D\) everywhere, then \(C\subset D\). In particular
 \(h_C=h_D\) implies \(C=D\).

 \emph{Proof.} Suppose \(x\notin D\) and let \(y\in D\) be a point of
 \(D\) nearest to \(x\). For \(z\in D\) and \(t\in(0,1]\) the point
 \(y+t(z-y)\) lies in \(D\), so
 \(\|x-y\|^2\le\|x-y-t(z-y)\|^2
 =\|x-y\|^2-2t\langle x-y,z-y\rangle+t^2\|z-y\|^2\). Dividing by \(t\)
 and letting \(t\to0\) gives \(\langle x-y,z-y\rangle\le0\). Hence
 \(h_D(x-y)=\langle x-y,y\rangle
 <\langle x-y,y\rangle+\|x-y\|^2=\langle x-y,x\rangle\). If \(x\) were
 in \(C\) we would get \(h_C(x-y)\ge\langle x-y,x\rangle>h_D(x-y)\),
 contrary to hypothesis. So \(C\subset D\). \qed
\item \(h_{[-v/2,\,v/2]}(u)=\tfrac12|\langle v,u\rangle|\), because a
 linear function on a segment attains its maximum at an endpoint.
\end{enumerate}

\textbf{Polytope facts.} A \emph{convex polytope} is the convex hull of
finitely many points. For a full-dimensional polytope \(P\subset\R^d\)
and \(u\in\R^d\), the \emph{face exposed by} \(u\) is
\(P(u)=\{x\in P:\langle u,x\rangle=h_P(u)\}\); thus \(P(0)=P\), and
for \(u\ne0\) the face \(P(u)\) lies in a hyperplane, so
\(\dim P(u)\le d-1\). A \emph{facet} is a face of dimension \(d-1\); it
equals \(P(\nu)\) for its outward unit normal \(\nu\). We use the
standard facts that \(P\) has finitely many faces, that the boundary of
\(P\) is the union of its facets, and that two distinct facets meet in
a face of dimension at most \(d-2\).

\textbf{Volume facts.} \(|MC|=|\det M|\,|C|\) for linear \(M\);
\(|tC|=t^d|C|\) in \(\R^d\); \(|C\times D|=|C|\,|D|\).

\section{The facet formula}\label{sec:facets}

This lemma is the polytopal case of Cauchy's projection formula; see
\cite[Section~1]{Saroglou}. It is proved here in full because the
factor \(\tfrac12\) fixes all later constants. It is used in
Proposition~\ref{prop:product}, Lemma~\ref{lem:affine} and
Proposition~\ref{prop:simplex}.

\begin{lemma}\label{lem:facets}
Let \(P\subset\R^d\), \(d\ge2\), be a full-dimensional convex polytope.
For each facet \(F\), let \(s_F\) be its \((d-1)\)-dimensional volume
and \(\nu_F\) its outward unit normal. Then
\begin{equation}\label{eq:facet}
 h_{\Pi P}(u)=\frac12\sum_Fs_F|\langle\nu_F,u\rangle|,
 \qquad u\in\R^d,
\end{equation}
and therefore
\begin{equation}\label{eq:zonotope}
 \Pi P=\sum_F
 \left[-\frac{s_F\nu_F}{2},\frac{s_F\nu_F}{2}\right].
\end{equation}
\end{lemma}

\begin{proof}
\emph{Plan.} For a unit vector \(u\), project each facet to \(u^\perp\).
The projected facet has volume \(s_F|\langle\nu_F,u\rangle|\) (Step~1),
and almost every point of the shadow of \(P\) is covered by exactly two
projected facets (Step~2). Integrating gives \eqref{eq:facet} (Step~3).
Fix a unit vector \(u\) and write \(\pi=\proj_{u^\perp}\).

\step{Step 1 (Jacobian).} \emph{Claim: for every facet \(F\),
\(\vol_{d-1}(\pi F)=s_F|\langle\nu_F,u\rangle|\). If
\(\langle\nu_F,u\rangle=0\), then \(\pi F\) lies in an affine subspace
of \(u^\perp\) of dimension at most \(d-2\).}

Let \(v_1,\ldots,v_{d-1}\) be an orthonormal basis of \(\nu_F^\perp\),
the direction space of the hyperplane \(H\) of \(F\). The restriction
\(\pi|_H\colon H\to u^\perp\) is an affine map between
\((d-1)\)-dimensional spaces, so it multiplies \((d-1)\)-volume by the
\((d-1)\)-volume \(J\) of the parallelepiped spanned by
\(\pi v_1,\ldots,\pi v_{d-1}\). Since \(u\) is a unit vector orthogonal
to all \(\pi v_i\), \(J\) equals the \(d\)-volume of the parallelepiped
spanned by \(\pi v_1,\ldots,\pi v_{d-1},u\):
\[
 J=|\det(\pi v_1,\ldots,\pi v_{d-1},u)|
  =|\det(v_1,\ldots,v_{d-1},u)|,
\]
where the second equality holds because
\(\pi v_i=v_i-\langle v_i,u\rangle u\) differs from \(v_i\) by a
multiple of the last column. Expand
\(u=\langle u,\nu_F\rangle\nu_F+\sum_i\langle u,v_i\rangle v_i\) in the
last column. The terms with \(v_i\) repeat a column and vanish, so
\(J=|\langle u,\nu_F\rangle|\cdot|\det(v_1,\ldots,v_{d-1},\nu_F)|
=|\langle\nu_F,u\rangle|\), the last matrix being orthogonal. If
\(J\ne0\), \(\pi|_H\) is an affine bijection and
\(\vol_{d-1}(\pi F)=Js_F\). If \(J=0\), the linear part of \(\pi|_H\)
is singular, so \(\pi H\) is a proper affine subspace of \(u^\perp\)
and \(\vol_{d-1}(\pi F)=0=Js_F\).

\step{Step 2 (multiplicity two).} Let \(E\subset u^\perp\) be the union
of
\begin{enumerate}[label=(\roman*),leftmargin=3em,nosep]
\item the relative boundary of \(\pi P\) in \(u^\perp\);
\item the sets \(\pi G\) for all faces \(G\) of \(P\) with
 \(\dim G\le d-2\);
\item the sets \(\pi F\) for all facets \(F\) with
 \(\langle\nu_F,u\rangle=0\).
\end{enumerate}
\emph{\(E\) has \((d-1)\)-dimensional measure zero.} The boundary of a
convex set is a null set, which covers (i). The sets in (ii) and (iii)
are finitely many and each lies in an affine subspace of dimension at
most \(d-2\) (for (iii) by Step~1).

\emph{Claim: every \(y\in\pi P\setminus E\) lies in \(\pi F\) for
exactly two facets \(F\).} Let \(\ell=y+\R u\).

(a) \emph{\(\ell\) meets the interior of \(P\).} Otherwise the
separation theorem gives \(\xi\ne0\) and \(c\in\R\) with
\(\langle\xi,x\rangle\le c\le\langle\xi,z\rangle\) for all \(x\in P\),
\(z\in\ell\). Boundedness below on the line \(\ell\) forces
\(\langle\xi,u\rangle=0\), so \(\xi\in u^\perp\) and
\(\langle\xi,\pi x\rangle=\langle\xi,x\rangle\le c\le\langle\xi,y\rangle\)
for all \(x\in P\). Then \(y\) maximizes a nonzero linear functional of
\(u^\perp\) over \(\pi P\), so \(y\) is a relative boundary point of
\(\pi P\), contradicting (i).

(b) By (a) and convexity, \(\ell\cap P\) is a segment \([p^-,p^+]\)
with \(p^-\ne p^+\). Its endpoints lie in the boundary of \(P\) and all
its other points lie in the interior of \(P\). (If \(c\) is an interior
point and \(p\in P\), every point of \([c,p]\) other than \(p\) is
interior.) Thus \(\ell\) meets the boundary of \(P\) exactly in
\(\{p^-,p^+\}\).

(c) The boundary point \(p^-\) lies in some facet. It lies in no face
of dimension at most \(d-2\), since \(\pi p^-=y\) avoids the sets (ii).
As two distinct facets meet in such a face, \(p^-\) lies in exactly one
facet \(F^-\). Likewise \(p^+\) lies in exactly one facet \(F^+\).

(d) \(F^-\ne F^+\): a facet containing both \(p^-\) and \(p^+\) has a
hyperplane containing the direction \(p^+-p^-\), a nonzero multiple of
\(u\), so \(\langle\nu_F,u\rangle=0\) and \(y\in\pi F\subset E\) by
(iii), a contradiction.

(e) If \(y\in\pi F\) for a facet \(F\), then \(F\) contains a point of
\(\ell\), which is a boundary point of \(P\), hence \(p^-\) or \(p^+\)
by (b). So \(F\in\{F^-,F^+\}\) by (c). Conversely \(y\in\pi F^\pm\).
With (d), exactly two facets qualify.

\step{Step 3 (integration).} Write \(1_X\) for the indicator function
of a set \(X\). Each \(\pi F\subset\pi P\), so by Step~1, Step~2 and
the nullity of \(E\),
\[
 \sum_Fs_F|\langle\nu_F,u\rangle|
 =\sum_F\vol_{d-1}(\pi F)
 =\int_{u^\perp}\sum_F1_{\pi F}(y)\,dy
 =\int_{\pi P\setminus E}2\,dy
 =2\vol_{d-1}(\pi P).
\]
This is \eqref{eq:facet} for unit \(u\). Both sides of
\eqref{eq:facet} are positively homogeneous of degree one in \(u\) (the
left side by definition) and vanish at \(u=0\), so \eqref{eq:facet}
holds for all \(u\in\R^d\).

\step{Step 4 (zonotope).} By (S3) and (S1), the right side of
\eqref{eq:facet} is the support function of the Minkowski sum in
\eqref{eq:zonotope}. So a compact convex set with the support function
required of \(\Pi P\) exists, and by (S2) it is that sum.
\end{proof}

\section{The product identity}\label{sec:product}

This proposition is the mechanism of the counterexample. It is used in
the proofs of Theorem~\ref{thm:counterexample} and
Corollary~\ref{cor:exponential-excess}.

\begin{proposition}\label{prop:product}
Let \(A\subset\R^r\) and \(B\subset\R^s\) be full-dimensional convex
polytopes, where \(r,s\ge2\). In the orthogonal product space
\(\R^{r+s}=\R^r\times\R^s\),
\begin{equation}\label{eq:product-body}
 \Pi(A\times B)=(|B|\Pi A)\times(|A|\Pi B).
\end{equation}
In particular,
\begin{equation}\label{eq:multiplicative}
 R_{r+s}(A\times B)=R_r(A)R_s(B).
\end{equation}
\end{proposition}

\begin{proof}
\emph{Plan.} Identify the facets of \(A\times B\) with their areas and
normals, insert them into the facet formula \eqref{eq:facet}, and read
off a product. Then compare volumes.

\step{Step 1 (facets of the product).} \emph{Claim: the facets of
\(A\times B\) are exactly the sets \(F\times B\) with \(F\) a facet of
\(A\), and \(A\times G\) with \(G\) a facet of \(B\).}

For \((u,v)\in\R^r\times\R^s\) and \((x,y)\in A\times B\) we have
\(\langle(u,v),(x,y)\rangle=\langle u,x\rangle+\langle v,y\rangle\),
which is maximal exactly when both terms are. Hence the exposed faces
satisfy \((A\times B)(u,v)=A(u)\times B(v)\), and
\(\dim(A\times B)(u,v)=\dim A(u)+\dim B(v)\). A facet of \(A\times B\)
is \((A\times B)(u,v)\) for some \((u,v)\ne0\) and has dimension
\(r+s-1\). If \(u\ne0\) and \(v\ne0\), then \(\dim A(u)\le r-1\) and
\(\dim B(v)\le s-1\), so the dimension is at most \(r+s-2\); this case
is impossible. If \(v=0\ne u\), the face is \(A(u)\times B\), of
dimension \(\dim A(u)+s\), so \(A(u)\) is a facet \(F\) of \(A\). The
case \(u=0\ne v\) is symmetric. Conversely
\(F\times B=(A\times B)(\nu_F,0)\) has dimension \(r+s-1\), so it is a
facet, and likewise \(A\times G\).

\step{Step 2 (normals and areas).} The facet \(F\times B\) is exposed
by the unit vector \((\nu_F,0)\), which is therefore its outward unit
normal. It lies in the hyperplane
\(\mathrm{aff}(F)\times\R^s\), which is the orthogonal product of the
hyperplane \(\mathrm{aff}(F)\subset\R^r\) and \(\R^s\), so its
\((r+s-1)\)-volume is the product \(s_F|B|\). Symmetrically,
\(A\times G\) has outward unit normal \((0,\nu_G)\) and volume
\(|A|s_G\). The area-normal vectors are thus
\[
 (|B|s_F\nu_F,0)
 \quad\hbox{and}\quad
 (0,|A|s_G\nu_G).
\]

\step{Step 3 (support function).} By \eqref{eq:facet} applied to
\(A\times B\), then to \(A\) and to \(B\),
\[
 h_{\Pi(A\times B)}(u,v)
 =\frac12\sum_F|B|s_F|\langle\nu_F,u\rangle|
  +\frac12\sum_G|A|s_G|\langle\nu_G,v\rangle|
 =|B|h_{\Pi A}(u)+|A|h_{\Pi B}(v).
\]
By (S1) this is the support function of
\((|B|\Pi A)\times(|A|\Pi B)\), so \eqref{eq:product-body} follows from
(S2).

\step{Step 4 (volumes).} Since \(\Pi A\subset\R^r\) and
\(\Pi B\subset\R^s\), \eqref{eq:product-body} gives
\(|\Pi(A\times B)|=|B|^r|\Pi A|\cdot|A|^s|\Pi B|\). Also
\(|A\times B|=|A||B|\). Hence
\[
 R_{r+s}(A\times B)
 =\frac{|B|^r|A|^s|\Pi A||\Pi B|}
         {(|A||B|)^{r+s-1}}
 =\frac{|\Pi A|}{|A|^{r-1}}
   \cdot\frac{|\Pi B|}{|B|^{s-1}}
 =R_r(A)R_s(B),
\]
the exponents being \(s-(r+s-1)=-(r-1)\) for \(|A|\) and
\(r-(r+s-1)=-(s-1)\) for \(|B|\).
\end{proof}

\section{Affine invariance}\label{sec:affine}

To speak of ``the simplex value'' we need all simplices of a given
dimension to have the same \(R_d\). The following lemma supplies this.
It is used only in the last sentence of
Proposition~\ref{prop:simplex}, and there only for polytopes (Part~A
below).

\begin{lemma}\label{lem:affine}
Let \(P\subset\R^d\), \(d\ge2\), be a convex body. For every
invertible linear map \(L\colon\R^d\to\R^d\),
\[
 \Pi(LP)=|\det L|L^{-t}\Pi P,\qquad R_d(LP)=R_d(P),
\]
where \(L^{-t}\) denotes the inverse transpose. Translations also
preserve \(R_d\).
\end{lemma}

\begin{proof}
\emph{Plan.} Translations are immediate. Part~A proves the lemma for
polytopes by transforming each facet's area-normal vector and applying
\eqref{eq:zonotope}. Part~B extends it to arbitrary convex bodies by
inscribed polytopes, and also proves that \(\Pi P\) exists.

\step{Translations.} For \(t\in\R^d\),
\(\proj_{u^\perp}(P+t)\) is a translate of \(\proj_{u^\perp}P\), so the
two bodies have the same projection volumes, hence the same projection
body, and \(|P+t|=|P|\). So \(R_d(P+t)=R_d(P)\).

\medskip\noindent\textbf{Part A: \(P\) is a polytope.}

\step{A1 (normals).} \(L\) is a linear bijection, so the facets of
\(LP\) are the sets \(LF\), \(F\) a facet of \(P\). For \(x\in P\),
\(\langle L^{-t}\nu_F,Lx\rangle=\langle\nu_F,x\rangle\). Hence
\(L^{-t}\nu_F\) is maximized over \(LP\) exactly on \(LF\), and the
outward unit normal of \(LF\) is
\(\nu_{LF}=L^{-t}\nu_F/\|L^{-t}\nu_F\|\). This holds for either sign of
\(\det L\).

\step{A2 (areas).} \emph{Claim:
\(s_{LF}=|\det L|\,\|L^{-t}\nu_F\|\,s_F\).} Let
\(Z=F+[0,\nu_F]\), a prism of height one over \(F\); \(|Z|=s_F\). Its
image is \(LZ=LF+[0,L\nu_F]\), with \(|LZ|=|\det L|s_F\). The slices of
\(LZ\) parallel to \(LF\) are the translates \(LF+tL\nu_F\),
\(t\in[0,1]\), at distance
\(t\langle L\nu_F,\nu_{LF}\rangle
=t\langle\nu_F,\nu_F\rangle/\|L^{-t}\nu_F\|=t/\|L^{-t}\nu_F\|\) from
the hyperplane of \(LF\). By Fubini,
\(|LZ|=s_{LF}/\|L^{-t}\nu_F\|\). Equating the two expressions for
\(|LZ|\) gives the claim.

\step{A3 (covariance).} By A1 and A2,
\(s_{LF}\nu_{LF}=|\det L|\,L^{-t}(s_F\nu_F)\). Put
\(M=|\det L|L^{-t}\). A linear map sends segments to segments and
Minkowski sums to Minkowski sums, so \eqref{eq:zonotope} gives
\[
 \Pi(LP)=\sum_F\Bigl[-\tfrac12M(s_F\nu_F),\tfrac12M(s_F\nu_F)\Bigr]
 =M\,\Pi P .
\]

\step{A4 (invariance).} \(|\det M|=|\det L|^d|\det L|^{-1}
=|\det L|^{d-1}\), so \(|\Pi(LP)|=|\det L|^{d-1}|\Pi P|\). Also
\(|LP|^{d-1}=|\det L|^{d-1}|P|^{d-1}\). Dividing,
\(R_d(LP)=R_d(P)\).

\medskip\noindent\textbf{Part B: \(P\) is an arbitrary convex body.}
(Not needed for Theorem~\ref{thm:counterexample} or
Corollary~\ref{cor:exponential-excess}.)

\step{B1 (approximating polytopes).} By translation invariance we may
assume that \(0\) is an interior point of \(P\). Choose \(\rho>0\) with
\(P\subset\rho B_2^d\), and a polytope \(P_0\subset P\) containing a
ball \(aB_2^d\), \(a>0\) (for instance a small cube centered at \(0\)).
For each \(m\ge1\) choose a finite set \(N_m\subset P\) such that every
point of \(P\) is within \(\varepsilon_m:=1/m\) of \(N_m\)
(compactness), and let \(P_m\) be the convex hull of
\(P_0\cup N_1\cup\cdots\cup N_m\). Then the \(P_m\) are polytopes,
\(aB_2^d\subset P_m\subset P_{m+1}\subset P\), and
\[
 P_m\subset P\subset P_m+\varepsilon_mB_2^d
 \subset\tau_mP_m,
 \qquad \tau_m:=1+\varepsilon_m/a\longrightarrow1 .
\]
The last inclusion holds because
\(\varepsilon_mB_2^d=(\varepsilon_m/a)\,aB_2^d\subset(\varepsilon_m/a)P_m\)
and \(P_m+(\varepsilon_m/a)P_m=\tau_mP_m\) by convexity.

\step{B2 (projection volumes).} For unit \(u\) put
\(g(u)=\vol_{d-1}(\proj_{u^\perp}P)\) and
\(g_m(u)=h_{\Pi P_m}(u)=\vol_{d-1}(\proj_{u^\perp}P_m)\). Projecting
\(P_m\subset P\subset\tau_mP_m\) and taking \((d-1)\)-volumes gives
\begin{equation}\label{eq:sandwich}
 g_m(u)\le g(u)\le\tau_m^{d-1}g_m(u).
\end{equation}
Since \(\proj_{u^\perp}P_m\) lies in a \((d-1)\)-ball of radius
\(\rho\), \(g_m(u)\le\kappa_{d-1}\rho^{d-1}\), and therefore
\[
 0\le g(u)-g_m(u)
 \le (\tau_m^{d-1}-1)\kappa_{d-1}\rho^{d-1}.
\]
Thus \(g_m\to g\) uniformly on \(S^{d-1}\). The sequence \(g_m\) is
nondecreasing, because \(P_m\subset P_{m+1}\).

\step{B3 (existence of \(\Pi P\)).} By (S2) and monotonicity of
\(g_m\), the bodies \(\Pi P_m\) increase. They lie in the ball
\(\kappa_{d-1}\rho^{d-1}B_2^d\) (as \(g_m\le\kappa_{d-1}\rho^{d-1}\))
and contain \(\kappa_{d-1}a^{d-1}B_2^d\) (as
\(g_m(u)\ge\vol_{d-1}(\proj_{u^\perp}aB_2^d)=\kappa_{d-1}a^{d-1}\)),
again by (S2). The union of an increasing family of convex sets is
convex, so the closure \(Z\) of \(\bigcup_m\Pi P_m\) is a convex body.
For unit \(u\),
\(h_Z(u)=\sup_mh_{\Pi P_m}(u)=\sup_mg_m(u)=g(u)\) by B2. Hence the
homogeneous extension of \(g\) is the support function of \(Z\); that
is, \(\Pi P\) exists and equals \(Z\).

\step{B4 (volume control).} By \eqref{eq:sandwich}, (S1) and (S2),
\[
 \Pi P_m\subset\Pi P\subset\tau_m^{d-1}\Pi P_m,
 \qquad
 |\Pi P_m|\le|\Pi P|\le\tau_m^{d(d-1)}|\Pi P_m|,
\]
and from B1, \(|P_m|\le|P|\le\tau_m^d|P_m|\). Consequently
\(|\Pi P_m|\to|\Pi P|\), \(|P_m|\to|P|\), and \(R_d(P_m)\to R_d(P)\).

\step{B5 (the same for \(LP\)).} Applying \(L\) to B1 gives polytopes
\(LP_m\), increasing, with \(LP_m\subset LP\subset\tau_mLP_m\),
\(LP\subset\|L\|\rho B_2^d\), and
\(LP_m\supset(a/\|L^{-1}\|)B_2^d\). Steps B2--B4 used only these
properties, so they hold with \(LP,LP_m\) in place of \(P,P_m\). In
particular \(h_{\Pi(LP_m)}\to h_{\Pi(LP)}\) pointwise and
\(R_d(LP_m)\to R_d(LP)\).

\step{B6 (limit).} By Part~A, \(\Pi(LP_m)=|\det L|L^{-t}\Pi P_m\), so
by (S1), for every \(u\),
\[
 h_{\Pi(LP_m)}(u)=|\det L|\,h_{\Pi P_m}(L^{-1}u).
\]
Letting \(m\to\infty\) (B2 for \(P\), B5 for \(LP\), and homogeneity)
gives \(h_{\Pi(LP)}(u)=|\det L|\,h_{\Pi P}(L^{-1}u)\), which by (S1)
is the support function of \(|\det L|L^{-t}\Pi P\). By (S2),
\(\Pi(LP)=|\det L|L^{-t}\Pi P\). Finally \(R_d(LP_m)=R_d(P_m)\) by
Part~A, and the limits in B4 and B5 give \(R_d(LP)=R_d(P)\).
\end{proof}

\section{The simplex value}\label{sec:simplex}

This proposition gives the number with which the product is compared.
It is used in the proofs of Theorem~\ref{thm:counterexample} (with
\(d=10\), and with \(d=20\) to interpret \(c_{20}\)) and of
Corollary~\ref{cor:exponential-excess}.

\begin{proposition}\label{prop:simplex}
For \(d\ge2\), the standard simplex
\(T_d=\conv(0,e_1,\ldots,e_d)\subset\R^d\) satisfies
\[
 |T_d|=\frac1{d!},\qquad
 |\Pi T_d|=\frac{d+1}{((d-1)!)^d},\qquad
 R_d(T_d)=c_d.
\]
Every full-dimensional \(d\)-simplex has the same value \(c_d\).
\end{proposition}

\begin{proof}
\emph{Plan.} List the \(d+1\) area-normal vectors of \(T_d\). By
\eqref{eq:zonotope}, \(\Pi T_d\) is a scaled translate of the unit cube
plus one segment, whose volume is found by slicing along that segment.
Write \(w=e_1+\cdots+e_d\) and
\(T_d=\{x:x_i\ge0\ \forall i,\ \sum_ix_i\le1\}\).

\step{Step 1 (volume).} \(T_d\) is the cone with apex \(e_d\) and height
\(1\) over the base \(T_d\cap\{x_d=0\}\), a copy of \(T_{d-1}\). Its
slice at height \(t\) is a copy of \((1-t)T_{d-1}\), so
\(|T_d|=|T_{d-1}|\int_0^1(1-t)^{d-1}dt=|T_{d-1}|/d\). With
\(|T_1|=1\) this gives \(|T_d|=1/d!\).

\step{Step 2 (facets).} A \(d\)-simplex has \(d+1\) facets. For \(T_d\)
they are:
\begin{itemize}[nosep,leftmargin=2em]
\item the \(d\) coordinate facets \(T_d\cap\{x_i=0\}\). Each is a copy
 of \(T_{d-1}\), with area \(1/(d-1)!\) by Step~1 and outward unit
 normal \(-e_i\). Its area-normal vector is \(-e_i/(d-1)!\).
\item the facet \(F_0=T_d\cap\{\sum_ix_i=1\}\), with outward unit
 normal \(w/\sqrt d\). Dropping the last coordinate is the orthogonal
 projection onto \(e_d^\perp\). It maps \(F_0\) bijectively onto
 \(\{x_1,\ldots,x_{d-1}\ge0,\ \sum_{i<d}x_i\le1\}\), a copy of
 \(T_{d-1}\). By Step~1 of the proof of Lemma~\ref{lem:facets} this
 projection scales area by
 \(|\langle w/\sqrt d,e_d\rangle|=1/\sqrt d\). Hence
 \(s_{F_0}/\sqrt d=1/(d-1)!\), i.e.\ \(s_{F_0}=\sqrt d/(d-1)!\), and the
 area-normal vector is \(s_{F_0}\,w/\sqrt d=w/(d-1)!\).
\end{itemize}

\step{Step 3 (shape of \(\Pi T_d\)).} By \eqref{eq:zonotope}, and since
\([-v/2,v/2]\) is unchanged when \(v\) is replaced by \(-v\),
\[
 \Pi T_d=\frac1{(d-1)!}\Bigl(\sum_{i=1}^d
 \bigl[-\tfrac{e_i}2,\tfrac{e_i}2\bigr]
 +\bigl[-\tfrac w2,\tfrac w2\bigr]\Bigr).
\]
Each centered segment \([-v/2,v/2]\) is the translate \(-v/2+[0,v]\),
and \(\sum_i[0,e_i]=[0,1]^d\). Therefore \(\Pi T_d\) is a translate (by
\(-w/(d-1)!\)) of
\[
 \frac{[0,1]^d+[0,w]}{(d-1)!}.
\]

\step{Step 4 (volume of the cube plus a segment).} Let \(Q=[0,1]^d\)
and \(\omega=w/\|w\|\). \emph{Claim: \(|Q+[0,w]|=d+1\).} For
\(y\in w^\perp\), the fiber \(\{t\in\R:y+t\omega\in Q\}\) is empty if
\(y\notin\proj_{w^\perp}Q\), and otherwise a compact interval
\([\alpha(y),\beta(y)]\). A point of \(Q+[0,w]\) is \(x+\sigma w\) with
\(x\in Q\), \(\sigma\in[0,1]\); it has the same projection to
\(w^\perp\) as \(x\), and its \(\omega\)-coordinate is that of \(x\)
plus \(\sigma\|w\|\). So the fiber of \(Q+[0,w]\) over \(y\) is empty
for \(y\notin\proj_{w^\perp}Q\) and equals
\([\alpha(y),\beta(y)+\|w\|]\) otherwise. By Fubini,
\[
 |Q+[0,w]|
 =\int_{\proj_{w^\perp}Q}\bigl(\beta(y)-\alpha(y)+\|w\|\bigr)dy
 =|Q|+\|w\|\vol_{d-1}(\proj_{w^\perp}Q)
 =1+h_{\Pi Q}(w),
\]
the last step by the definition of \(h_{\Pi Q}\) and its homogeneity.
The cube has \(2d\) facets, each of area one, with normals \(\pm e_i\),
and \(|\langle\pm e_i,w\rangle|=1\). So \eqref{eq:facet} gives
\(h_{\Pi Q}(w)=\tfrac12\cdot2d=d\), and the claim follows.
Figure~\ref{fig:extrusion} shows the case \(d=2\).

\step{Step 5 (the value).} By Steps~3 and~4,
\(|\Pi T_d|=(d+1)/((d-1)!)^d\). With Step~1 and \(d!=d\,(d-1)!\),
\[
 R_d(T_d)=\frac{(d+1)(d!)^{d-1}}{((d-1)!)^d}
 =\frac{(d+1)d^{d-1}}{(d-1)!}
 =\frac{(d+1)d^d}{d!}=c_d.
\]
\emph{Sanity check, \(d=2\):} \(|\Pi T_2|=3\), \(|T_2|=\tfrac12\),
\(R_2(T_2)=6=c_2\).

\step{Step 6 (all simplices).} A full-dimensional simplex
\(\conv(p_0,\ldots,p_d)\) equals \(p_0+LT_d\), where \(L\) is the linear
map with \(Le_i=p_i-p_0\); \(L\) is invertible because the \(p_i\) are
affinely independent. By Lemma~\ref{lem:affine} (Part~A and
translations), its value is \(R_d(T_d)=c_d\).
\end{proof}

\begin{figure}[tb]
\centering
\begin{tikzpicture}[scale=1.5, line join=round]
 \fill[black!7] (0,0) rectangle (1,1);
 \fill[blue!13] (1,0)--(2,1)--(2,2)--(1,1)--cycle;
 \fill[yellow!23] (0,1)--(1,1)--(2,2)--(1,2)--cycle;
 \draw[thick] (0,0)--(1,0)--(2,1)--(2,2)--(1,2)--(0,1)--cycle;
 \draw[thin] (0,1)--(1,1)--(1,0);
 \draw[thin] (1,1)--(2,2);
 \node at (.5,.5) {\(Q\)};
 \node at (1.5,1) {\(1\)};
 \node at (1,1.5) {\(1\)};
 \node[below left] at (0,0) {\(0\)};
 \node[above right] at (2,2) {\((2,2)\)};
 \draw[->,thick] (2.6,.3)--(3.6,1.3);
 \node[right] at (3.15,.65) {\(w=(1,1)\)};
\end{tikzpicture}
\caption{The planar extrusion \(Q+[0,(1,1)]\), with
\(Q=[0,1]^2\), consists of the unit square and two parallelograms of
area one. Its area is \(3\). The fiber argument in
Proposition~\ref{prop:simplex} proves the corresponding volume \(d+1\)
in every dimension.}
\label{fig:extrusion}
\end{figure}

\section{The counterexample and exponential excess}\label{sec:counterexample}

\subsection{Proof of the main theorem}

\begin{proof}[Proof of Theorem~\ref{thm:counterexample}]
\step{Step 1.} \(T_{10}\) is a full-dimensional convex polytope in
\(\R^{10}\), so \(K=T_{10}\times T_{10}\) is a full-dimensional convex
polytope in \(\R^{20}\), in particular a convex body.

\step{Step 2.} Proposition~\ref{prop:product} with \(r=s=10\),
\(A=B=T_{10}\), and then Proposition~\ref{prop:simplex} with \(d=10\),
give
\[
 R_{20}(K)=R_{10}(T_{10})^2=c_{10}^2
         =\left(\frac{11\cdot10^{10}}{10!}\right)^2.
\]

\step{Step 3.} Since \(c_{20}=21\cdot20^{20}/20!\),
\(10^{20}/20^{20}=2^{-20}\) and \(20!/(10!)^2=\binom{20}{10}\),
\[
 \frac{R_{20}(K)}{c_{20}}
 =\frac{121\cdot10^{20}\cdot20!}
        {(10!)^2\cdot21\cdot20^{20}}
 =\frac{121\binom{20}{10}}{21\cdot2^{20}}.
\]

\step{Step 4.} \(\binom{20}{10}=184\,756\) and \(2^{20}=1\,048\,576\),
so the numerator is \(121\cdot184\,756=22\,355\,476\) and the
denominator is \(21\cdot1\,048\,576=22\,020\,096\). The exact integer
certificate is
\[
 121\cdot184\,756-21\cdot1\,048\,576=335\,380>0.
\]
Hence \(|\Pi K|>c_{20}|K|^{19}\), a strict violation of
\eqref{eq:proposed-bound}. By Proposition~\ref{prop:simplex} with
\(d=20\), \(c_{20}\) is the value of every \(20\)-dimensional simplex.
\end{proof}

\subsection{The general product test}

For \(r,s\ge2\), Propositions~\ref{prop:product} and~\ref{prop:simplex}
give \(R_{r+s}(T_r\times T_s)=c_rc_s\). Dividing by
\(c_{r+s}=(r+s+1)(r+s)^{r+s}/(r+s)!\) and using
\((r+s)!/(r!\,s!)=\binom{r+s}{r}\),
\[
 \frac{R_{r+s}(T_r\times T_s)}{c_{r+s}}
 =\frac{(r+1)(s+1)}{r+s+1}
   \binom{r+s}{r}\frac{r^rs^s}{(r+s)^{r+s}}.
\]
Thus any universal upper constant in dimension \(r+s\) must be at least
\(c_rc_s\). \emph{Sanity checks:} for \(r=s=2\) the ratio is
\(\tfrac95\cdot6\cdot\tfrac1{16}=\tfrac{27}{40}<1\), and for \(r=s=9\)
it is \(\tfrac{100}{19}\cdot\tfrac{48\,620}{262\,144}<1\), so these
products do not violate \eqref{eq:proposed-bound}. For \(r=s=10\) the
formula reproduces Theorem~\ref{thm:counterexample}.

\subsection{Exponential excess in large dimensions}

\begin{corollary}\label{cor:exponential-excess}
There exist a constant \(\lambda>1\) and an integer \(N\ge20\) such that, for every
integer \(n\ge N\), some full-dimensional convex polytope
\(K_n\subset\R^n\) satisfies
\[
 \frac{R_n(K_n)}{R_n(T_n)}\ge\lambda^n.
\]
The polytopes \(K_n\) may be chosen as Cartesian products of simplices.
\end{corollary}

\begin{proof}
\emph{Plan.} Take \(k\) factors \(T_{10}\) and one further simplex of
dimension between \(10\) and \(19\) to reach dimension \(n\). The value
is \(c_{10}^k\) times a bounded factor. Bound \(c_n\) above by
\((n+1)\beta^n\) with \(\beta=11/4\), and check that
\(c_{10}>\beta^{10}\). Then each ten-dimensional factor gains a fixed
factor larger than one.

\step{Step 1 (the bodies).} Every integer \(n\ge20\) is uniquely
\(n=10k+q\) with integers \(k\ge1\) and \(q\in\{10,\ldots,19\}\)
(divide \(n-10\) by \(10\)). Put
\[
 K_n=T_{10}^{\,k}\times T_q,
\]
where \(T_{10}^{\,k}\) is the \(k\)-fold Cartesian product, in the
orthogonal coordinate product space.

\step{Step 2 (its value).} \emph{Claim: \(R_n(K_n)=c_{10}^kc_q\).} By
induction on \(k\), \(T_{10}^{\,k}\) is a full-dimensional convex
polytope with \(R_{10k}(T_{10}^{\,k})=c_{10}^k\). The case \(k=1\) is
Proposition~\ref{prop:simplex}, and the step from \(k\) to \(k+1\) is
Proposition~\ref{prop:product} applied to
\(T_{10}^{\,k}\times T_{10}\). One more application of
Proposition~\ref{prop:product}, to \(T_{10}^{\,k}\times T_q\), together
with Proposition~\ref{prop:simplex} for \(T_q\), proves the claim. All
factor dimensions are at least \(2\), as required.

\step{Step 3 (\(e<\beta:=11/4\)).} In \(e=\sum_{j\ge0}1/j!\) the terms
with \(j\le2\) sum to \(\tfrac52\). For \(j\ge3\) the ratio of
consecutive terms is \(1/(j+1)\le\tfrac14<\tfrac13\), so
\[
 e=\frac52+\sum_{j=3}^{\infty}\frac1{j!}
 <\frac52+\frac16\sum_{j=0}^{\infty}3^{-j}
 =\frac52+\frac16\cdot\frac32
 =\beta .
\]

\step{Step 4 (upper bound for \(c_n\)).} The term \(n^n/n!\) of the
series \(e^n=\sum_jn^j/j!\) is strictly smaller than the full sum of
positive terms. By Step~3,
\[
 c_n=(n+1)\frac{n^n}{n!}<(n+1)e^n<(n+1)\beta^n\qquad(n\ge2).
\]

\step{Step 5 (the ten-dimensional block).} Since
\(10!=2^8\cdot3^4\cdot5^2\cdot7\) and \(10^{10}=2^{10}5^{10}\),
\[
 c_{10}=\frac{11\cdot10^{10}}{10!}
 =\frac{11\cdot2^2\cdot5^8}{3^4\cdot7}
 =\frac{17\,187\,500}{567}.
\]
With \(\beta^{-10}=4^{10}/11^{10}\),
\[
 \Theta:=\frac{c_{10}}{\beta^{10}}
 =\frac{2^{22}\cdot5^8}{3^4\cdot7\cdot11^9}
 =\frac{1\,638\,400\,000\,000}{1\,336\,956\,340\,797}>1 ,
\]
using \(2^{22}5^8=2^{14}\cdot10^8\) and
\(567\cdot11^9=567\cdot2\,357\,947\,691\).

\step{Step 6 (comparison).} Fix \(\lambda\) with
\(1<\lambda<\Theta^{1/10}\); this is possible since \(\Theta>1\). By
Proposition~\ref{prop:simplex}, \(R_n(T_n)=c_n\). Using Step~2, then
Step~4, then \(n=10k+q\),
\[
 \frac{R_n(K_n)}{R_n(T_n)\lambda^n}
 =\frac{c_{10}^kc_q}{c_n\lambda^n}
 >\frac{c_{10}^kc_q}{(n+1)\beta^{10k+q}\lambda^{10k+q}}
 =c_q\beta^{-q}\lambda^{-q}\,
   \frac{(\Theta/\lambda^{10})^k}{10k+q+1}.
\]

\step{Step 7 (choice of \(N\)).} For each fixed
\(q\in\{10,\ldots,19\}\), the right side tends to infinity as
\(k\to\infty\), because \(\Theta/\lambda^{10}>1\) and an exponential
outgrows a linear function. So there is \(k_q\ge1\) such that the right
side is at least \(1\) for all \(k\ge k_q\). Let
\(N=\max_q(10k_q+q)\), which is at least \(20\). If \(n\ge N\) and
\(n=10k+q\), then \(10k+q\ge10k_q+q\), so \(k\ge k_q\) and
\(R_n(K_n)/R_n(T_n)>\lambda^n\).
\end{proof}

\section{Verification checklist}

A sceptical reader should check these load-bearing points first.

\begin{enumerate}[leftmargin=2em]
\item \textbf{The factor \(\tfrac12\) in the facet formula.} Almost
 every point of the shadow is covered by exactly two projected facets
 (Lemma~\ref{lem:facets}, Step~2), and a projected facet has volume
 \(s_F|\langle\nu_F,u\rangle|\) (Step~1).
\item \textbf{Facets of a product and the exponent count.} The facets
 of \(A\times B\) are \(F\times B\) and \(A\times G\), with areas
 \(s_F|B|\) and \(|A|s_G\) (Proposition~\ref{prop:product}, Steps~1--2),
 and the powers of \(|A|,|B|\) cancel as claimed (Step~4).
\item \textbf{The slanted facet of the simplex.} Its area-normal vector
 is \(w/(d-1)!\) (Proposition~\ref{prop:simplex}, Step~2), which makes
 \(\Pi T_d\) a translate of \(([0,1]^d+[0,w])/(d-1)!\) (Step~3).
\item \textbf{The volume \(|[0,1]^d+[0,w]|=d+1\).} Fiber argument and
 \(h_{\Pi Q}(w)=d\) (Proposition~\ref{prop:simplex}, Step~4); compare
 Figure~\ref{fig:extrusion} for \(d=2\).
\item \textbf{The integer inequality.}
 \(121\cdot184\,756=22\,355\,476>22\,020\,096=21\cdot2^{20}\)
 (proof of Theorem~\ref{thm:counterexample}, Steps~3--4).
\item \textbf{For ``all simplices have value \(c_d\)''.} The area
 transformation \(s_{LF}=|\det L|\,\|L^{-t}\nu_F\|\,s_F\)
 (Lemma~\ref{lem:affine}, A2). Part~B of that lemma is not needed for
 the theorem.
\item \textbf{For the corollary.} \(e<11/4\) (Step~3),
 \(\Theta=c_{10}/(11/4)^{10}>1\) (Step~5), and the algebra in Step~6.
\end{enumerate}

\begin{thebibliography}{7}

\bibitem{Brannen}
N.~S.~Brannen,
\emph{Volumes of projection bodies},
Mathematika \textbf{43} (1996), no.~2, 255--264.
\href{https://doi.org/10.1112/S002557930001175X}{doi:10.1112/S002557930001175X}.

\bibitem{ChenEtAl}
S.~Chen, Y.~Feng, Y.~Li, D.~Xi, and L.~Xu,
\emph{Petty's conjectured projection inequality in dimension three},
preprint, 2026, MathSciDoc:2608.23001.
\href{https://archive.ymsc.tsinghua.edu.cn/pacm_download/741/12780-pettyinequality.pdf}{Full text}.

\bibitem{FengEtAl}
Y.~Feng, S.~Hu, W.~Liu, and L.~Xu,
\emph{On the reverse projection inequality},
preprint, 2026, MathSciDoc:2608.23002.
\href{https://archive.ymsc.tsinghua.edu.cn/pacm_download/743/12781--2026.8.26.pdf}{Full text}.

\bibitem{Henk}
M.~Henk,
\emph{Note on projection bodies of zonotopes with \(n+1\) generators},
Acta Mathematica Scientia \textbf{45} (2025), 96--103.
\href{https://doi.org/10.1007/s10473-025-0107-9}{doi:10.1007/s10473-025-0107-9}.

\bibitem{Ludwig}
M.~Ludwig,
\emph{Projection bodies and valuations},
Advances in Mathematics \textbf{172} (2002), no.~2, 158--168.
\href{https://doi.org/10.1016/S0001-8708(02)00021-X}{doi:10.1016/S0001-8708(02)00021-X}.

\bibitem{Petty}
C.~M.~Petty,
\emph{Projection bodies},
in \emph{Proceedings of the Colloquium on Convexity (Copenhagen, 1965)},
K\o benhavns Universitet Matematiske Institut, Copenhagen, 1967,
pp.~234--241.

\bibitem{Saroglou}
C.~Saroglou,
\emph{On the shape of a convex body with respect to its second projection body},
Advances in Applied Mathematics \textbf{67} (2015), 55--74.
\href{https://doi.org/10.1016/j.aam.2015.03.004}{doi:10.1016/j.aam.2015.03.004}.
Preprint:
\href{https://arxiv.org/abs/1409.4347v2}{arXiv:1409.4347v2}.

\end{thebibliography}

\end{document}
